%% kapa: exam=1 round=35 year=2024 period=1 form=A subject=real_estate subject_ko=부동산학원론 seq=24 number=104
%% source: Q-Net 감정평가사 기출문제 (cst00309) artlSeq=5213229 file=raw/1008/5213229/1교시 문제지 원본.pdf
%% conversion: kordoc 4.24.51 -> tools/gen1.py (auto); review=proofread-sonnet; flags=none
\begin{question}{104}
부동산 증권에 관한 설명으로 옳은 것을 모두 고른 것은?
\begin{framedtext}
ㄱ. MPTS(Mortgage Pass-Through Securities)는 채권을 표시하는 증권으로 원리금수취권과 주택저당에 대한 채권을 모두 투자자에게 이전하는 증권이다. \par
ㄴ. MBB(Mortgage-Backed Bond)는 모기지 풀(Pool)에서 발생하는 현금흐름으로 채권의 원리금이 지급되고, 모기지 풀의 현금흐름으로 채권의 원리금지급이 안될 경우 발행자가 초과부담을 제공하는 채권이다. \par
ㄷ. CMO(Collateralized Mortgage Obligation)는 원금과 조기상환대금을 받아갈 순서를 정한 증권으로 증권별로 만기가 일치하도록 만든 자동이체형 증권이다. \par
ㄹ. MPTB(Mortgage Pay-Through Bond)는 채권으로 발행자의 대차대조표에 부채로 표시된다. \par
ㅁ. 금융기관은 MBS(Mortgage-Backed Securities)를 통해 자기자본비율(BIS)을 높일 수 있다. \par
\end{framedtext}
\begin{choices}
  \item ㄱ, ㄴ, ㄷ
  \item ㄱ, ㄴ, ㄹ
  \item ㄱ, ㄷ, ㅁ
  \item ㄴ, ㄹ, ㅁ
  \item ㄷ, ㄹ, ㅁ
\end{choices}
\end{question}
\answer{4}
